\section{Discussion}

\noindent\textbf{Reliability of the Corpus.}
One concern with our corpus is that some of the 12,707 PRs may not fix a bug, or that we may have assigned the wrong symptom, faulty layer, or repair operation.
% The labels attached to PRs on GitHub cannot settle this issue: a bug-related label may be missing or applied to a change that is not a bug fix.
We addressed this concern during screening: two experts independently examined the description, discussion, linked issue when available, and complete code change for each of the 50,121 candidate PRs. We retained a PR only if it was merged and the evidence showed that it corrected unintended behavior. The experts developed the classification categories from an initial 1,000 PRs, then independently labeled the retained fixes and resolved disagreements by discussion. Each final record keeps the PR link and the code change used to make the decision. With independent review, disagreement resolution, and retained PR evidence, we made every effort to include only genuine bug fixes and to label them accurately.

\noindent\textbf{Contribution of the Code-Centered Analysis.}
Our contribution lies in answering code-level maintenance questions that earlier studies of inference-engine bugs left open. Liu et al.~\cite{liu2026firstlook} already analyzed symptoms, root causes, fix strategies, fix effort, and temporal evolution. Those analyses do not show which files and functions repeatedly participate in repairs or whether complete fixes reach more files over time. We study a number of bug-fixing PRs and link each observed symptom, faulty architectural layer, and main repair operation to the files and functions in its complete patch. Using the same PR as the unit of analysis lets us examine where repair work accumulates and how its scope changes.

Our PR-level analysis shows that repairs recur in a small part of the codebase and that more fixes now touch several files. Across the six engines, the top 20\% of files account for 79.6\% of weighted bug-fix associations and the top 20\% of functions for 68.4\% (Fig.~\ref{fig:file_function_concentration}). From January 2023 to March 2026, the share of fixes modifying two to five files rises from 31.3\% to 41.8\% (Fig.~\ref{fig:monthly-file-bucket-trend}). Repair operations also differ by architectural layer: Logic Correction and Adaptation dominate most layers, while Configuration is especially common in Build/Release. These findings give maintainers evidence for prioritizing frequently repaired code during review and selecting regression tests that cover the complete fix.

The evidence serves different maintenance tasks. A characteristic symptom can point to a specialized boundary layer, whereas broad symptoms from Scheduler/Runtime or Compiler/Optimization require wider tracing. Within a candidate layer, file and function histories show recurring repair sites, and operation counts show common changes. The growth of two-to-five-file fixes also cautions against stopping review or regression testing at the first faulty entity. We have not evaluated this sequence as a diagnostic method; it is one way to interpret the observed distributions.

\noindent\textbf{Threats to Validity.}
Internal validity depends on our judgments about which PRs fix bugs and how those fixes are classified. A PR with limited explanation or several intertwined changes may receive the wrong symptom, faulty-layer, or repair-operation label; enough such errors would change the reported distributions and associations. We checked these decisions as carefully as the repository evidence permits: two experts independently reviewed the PR discussions and complete patches, resolved disagreements, and retained the evidence for each decision. Judgment cannot be eliminated, however. We record one primary symptom and repair operation per PR, so secondary effects may be missed, and PR counts do not necessarily equal counts of distinct defects.

External validity is limited by the projects and fixes we can observe. The corpus contains merged, public fixes from six open-source engines; undiscovered, unresolved, privately fixed, and externally fixed bugs are absent. vLLM and llama.cpp contribute 70.2\% of the PRs, so pooled percentages may differ in other engines or in a more balanced set of projects. We studied engines with different deployment settings and examined symptom profiles and repair concentration within each project. For the Hardware Adaptation trend, we also checked that the increase persists when project weights are held fixed. These checks show that both patterns extend beyond one project, although their exact percentages and slopes may vary across engines.
